
of @@M0@@, which follows from the geometric connectedness of @@M1@@.

From the homomorphism (7.12.1), using 7.4, one obtains a homomorphism
\sgaStableTarget{sga6.e10.eq.7-12-2}@@M2@@

When @@M3@@ and @@M4@@ are moreover assumed regular, (7.12.1) can also be interpreted as a ring homomorphism
\sgaStableTarget{sga6.e10.eq.7-12-3}@@M5@@
which, when @@M6@@ is moreover separated, is a homomorphism of @@M7@@-rings, and consequently defines a homomorphism
\sgaStableTarget{sga6.e10.eq.7-12-4}@@M8@@

\sgaStableTarget{sga6.e10.para.7-12-5}\textbf{7.12.5.} Let @@M9@@ be a smooth proper algebraic scheme over a field @@M10@@. One checks easily that an element @@M11@@ of @@M12@@, for which there is an integer @@M13@@ with @@M14@@ algebraically equivalent to zero, is numerically equivalent to zero. It is plausible that the converse is always true, at least when @@M15@@ is projective. In this question, one can plainly suppose that @@M16@@ is algebraically closed, and replace @@M17@@ by @@M18@@ or @@M19@@ without changing the conjecture. One also obtains an equivalent conjecture by working with the Chow ring @@M20@@ (cf. XIV 4.1), obtaining the following more or less classical conjecture: if @@M21@@ is an algebraic cycle numerically equivalent to zero on @@M22@@, there is an integer @@M23@@ with @@M24@@ algebraically equivalent to zero. In the case where this conjecture is correct for @@M25@@, the compatibility of the homomorphism (7.3.1) with algebraic equivalence would imply that this specialization homomorphism (when @@M26@@ is a projective smooth morphism) is also compatible with numerical equivalence, as conjectured in 7.9.
